\subsection{Multiplication as two-stage normalization: convolution and normal form}\label{subsec:emergent-arithmetic-stable-mul}

\begin{definition}[Stable multiplication]\label{def:stable-mul}
For \(c,d\in X_\infty\), define
\[
c\otimes d:=\mathrm{Val}^{-1}\bigl(\mathrm{Val}(c)\cdot\mathrm{Val}(d)\bigr)\in X_\infty.
\]
\leanverified{stableMul}
\end{definition}

\begin{proposition}[Fibonacci-kernel convolution followed by normal form]\label{prop:mul-convolution-normal-form}
For \(i,j\ge 1\), let
\[
K_{ij}:=Z(F_{i+1}F_{j+1})\in X_\infty,
\qquad
K_{ij,k}:=(K_{ij})_k.
\]
For \(c,d\in X_\infty\), define the finite raw convolution configuration
\[
\mathcal{C}(c,d):=
\sum_{i,j\ge 1}c_i d_j K_{ij}
=\sum_{k\ge 1}\left(\sum_{i,j\ge 1}c_i d_jK_{ij,k}\right)e_k
\in\mathcal{D}.
\]
Then
\[
c\otimes d=\Fold_\infty\bigl(\mathcal{C}(c,d)\bigr).
\]
Consequently stable multiplication admits a digit-level construction: form the finite Fibonacci-kernel convolution of the two normal words and then apply the finite local normalization procedure of Proposition \ref{prop:finite-local-fold-normalization}.
\end{proposition}

\begin{proof}
The sums defining \(\mathcal{C}(c,d)\) are finite because \(c\) and \(d\) have finite support, and each \(K_{ij}=Z(F_{i+1}F_{j+1})\) has finite support by Theorem \ref{thm:zeckendorf-unique}. Thus \(\mathcal{C}(c,d)\in\mathcal{D}\). By construction and additivity of \(\mathrm{Val}\) on \(\mathcal{D}\),
\[
\mathrm{Val}\bigl(\mathcal{C}(c,d)\bigr)
=\sum_{i,j\ge 1}c_i d_j\mathrm{Val}(K_{ij})
=\sum_{i,j\ge 1}c_i d_jF_{i+1}F_{j+1}.
\]
Factoring the finite double sum gives
\[
\mathrm{Val}\bigl(\mathcal{C}(c,d)\bigr)
=\left(\sum_{i\ge 1}c_iF_{i+1}\right)
 \left(\sum_{j\ge 1}d_jF_{j+1}\right)
=\mathrm{Val}(c)\mathrm{Val}(d).
\]
Therefore
\[
\Fold_\infty\bigl(\mathcal{C}(c,d)\bigr)
=Z\bigl(\mathrm{Val}(\mathcal{C}(c,d))\bigr)
=Z\bigl(\mathrm{Val}(c)\mathrm{Val}(d)\bigr)
=c\otimes d,
\]
where the last equality is Definition \ref{def:stable-mul}. Proposition \ref{prop:finite-local-fold-normalization} supplies the promised terminating local normalization from the raw convolution configuration to this normal form.
\end{proof}

\begin{theorem}[Fixed-resolution multiplication is finite-state, not a uniform online multiplier]\label{thm:finite-mul-fixed-resolution-fst}
Fix \(m\ge 1\) and set \(M_m:=F_{m+2}\). There is a finite transducer depending only on \(m\) which, on input two words \(c,d\in X_m\), outputs
\[
c\boxtimes_m d=Z_m\bigl(V_m(c)V_m(d)\bmod M_m\bigr).
\]
Equivalently, at each fixed resolution multiplication descends through the finite residue fold of Theorem \ref{thm:finite-resolution-mod}. This theorem is deliberately fixed-resolution: it supplies the finite-state realization needed for the finite algebra in the paper, but it is not a bounded-delay online multiplication theorem uniform in \(m\).
\end{theorem}

\begin{proof}
At fixed \(m\), the input language \(X_m\times X_m\) is finite. A concrete transducer may be described by its state register
\[
(r,s,u,v)\in \bigl(\ZZ/(M_m\ZZ)\bigr)^4,
\]
where after reading the first \(q\) low-to-high digit pairs the registers contain
\[
r\equiv \sum_{k=1}^{q}c_kF_{k+1},
\qquad
s\equiv \sum_{k=1}^{q}d_kF_{k+1},
\qquad
u\equiv F_{q+1},
\qquad
v\equiv F_{q+2}
\pmod {M_m}.
\]
The initial state is \((0,0,F_1,F_2)=(0,0,1,1)\bmod M_m\). On reading \((c_{q+1},d_{q+1})\), the machine updates
\[
r\leftarrow r+c_{q+1}v,
\qquad
s\leftarrow s+d_{q+1}v,
\qquad
(u,v)\leftarrow(v,u+v),
\]
modulo \(M_m\). After the last of the \(m\) digit pairs has been read, the registers contain \(r=V_m(c)\bmod M_m\) and \(s=V_m(d)\bmod M_m\). The machine then forms the residue \(rs\bmod M_m\) in its finite control and uses the fixed lookup table supplied by Lemma \ref{lem:zeckendorf-range-bijection},
\[
\{0,1,\dots,M_m-1\}\longrightarrow X_m,
\qquad n\longmapsto Z_m(n),
\]
to emit \(Z_m(rs\bmod M_m)\). By Definition \ref{def:finite-resolution-arithmetic}, this output is precisely \(c\boxtimes_m d\). All state sets and lookup tables are finite because \(m\) is fixed.
\end{proof}

\begin{theorem}[Definitional isomorphism for multiplication]\label{thm:mul-definitional}
\((X_\infty,\oplus,\otimes)\) is a commutative semiring, and \(\mathrm{Val}\) is a semiring isomorphism from \((X_\infty,\oplus,\otimes)\) onto \((\NN,+,\times)\).
\leanverified{paper\_stable\_commutative\_ring}
\leanverified{paper\_stableValue\_pow}
\leanverified{paper\_totient\_fib\_9\_10}
\leanverified{paper\_seven\_mul\_three\_zero\_X6}
\end{theorem}

\begin{proof}
By Definitions \ref{def:stable-add} and \ref{def:stable-mul}, the map \(\mathrm{Val}\) preserves both addition and multiplication. Since \(\mathrm{Val}\) is bijective, commutativity, associativity, and distributivity are transported from \((\NN,+,\times)\), giving the desired semiring isomorphism.
\end{proof}

Thus the natural numbers, addition, and multiplication are completely determined inside the stable address space by the syntactic constraints together with normalization to canonical form; the operations \(+\) and \(\times\) need not be inserted as primitive symbols.

\subsubsection{Multiplication generated by iterated addition}\label{subsec:mul-as-iterated-add}
\begin{theorem}[Multiplication generated by addition and iteration]\label{thm:mul-by-iterated-add}
Define, for \(c,d\in X_\infty\),
\[
c\odot d:=\underbrace{c\oplus c\oplus\cdots\oplus c}_{\mathrm{Val}(d)\text{ times}},
\qquad
\text{and take }0_X:=(0,0,\dots)\text{ if }\mathrm{Val}(d)=0.
\]
Then \(c\odot d=c\otimes d\) for all \(c,d\in X_\infty\).
\leanverified{iteratedStableAdd}
\leanverified{stableValue\_iteratedStableAdd}
\leanverified{iteratedStableAdd\_eq\_stableMul}
\end{theorem}

\begin{proof}
By Theorem \ref{thm:add-definitional}, the map \(\mathrm{Val}\) preserves \(\oplus\) and is bijective. Therefore
\[
\mathrm{Val}(c\odot d)=\mathrm{Val}(c)\mathrm{Val}(d)=\mathrm{Val}(c\otimes d),
\]
and hence \(c\odot d=c\otimes d\).
\end{proof}

\begin{corollary}[A first precise sense in which no new primitive is needed]\label{cor:mul-no-new-primitive}
The operation \(\otimes\) is uniquely determined by \(\oplus\) together with iterability of counting, and therefore need not be introduced as an additional primitive.
\end{corollary}

\subsubsection{Successor and folding: eliminating \texorpdfstring{$\oplus$}{oplus} as well}\label{subsec:succ-and-fold}
\begin{definition}[Successor and folding]\label{def:successor-fold}
Let \(0_X:=(0,0,\dots)\in X_\infty\), and let \(e_1\in\mathcal{D}\) be the first standard basis vector. Define
\[
S:X_\infty\to X_\infty,
\qquad
S(c):=\Fold_\infty(c+e_1).
\]
\leanverified{stableSucc}
\leanverified{stablePred}
\leanverified{stableValue\_stablePred}
\leanverified{stablePred\_stableSucc}
\leanverified{stableSucc\_stablePred}
\leanverified{stablePred\_injective}
\end{definition}

\begin{theorem}[Successor structure isomorphism]\label{thm:successor-structure}
The value map satisfies
\[
\mathrm{Val}(0_X)=0,
\qquad
\mathrm{Val}(S(c))=\mathrm{Val}(c)+1,
\]
and is bijective.
\leanverified{stableValue\_stableSucc}
\leanverified{stableSucc\_zero}
\leanverified{stableSucc\_injective}
\leanverified{stableSucc\_surjective}
\leanverified{stableSucc\_bijective}
\leanverified{stablePred\_surjective}
\leanverified{stablePred\_bijective}
\leanverified{stableSuccEquiv}
\end{theorem}

\begin{proof}
We have \(\mathrm{Val}(0_X)=0\). Moreover,
\[
\mathrm{Val}(S(c))
=\mathrm{Val}(\Fold_\infty(c+e_1))
=\mathrm{Val}(c+e_1)
=\mathrm{Val}(c)+\mathrm{Val}(e_1)
=\mathrm{Val}(c)+1,
\]
since \(\mathrm{Val}(e_1)=F_2=1\). Bijectivity again follows from uniqueness of Zeckendorf expansion.
\end{proof}

\begin{corollary}[Addition recovered from iterated successor]\label{cor:add-from-successor}
Define
\[
c\oplus_S d:=S^{\circ\mathrm{Val}(d)}(c).
\]
Then \(c\oplus_S d=c\oplus d\), and
\[
c\oplus d=\Fold_\infty(c+d).
\]
\leanverified{stableValue\_stableSucc\_iterate}
\leanverified{stableSucc\_iterate\_eq\_stableAdd}
\end{corollary}

\begin{proof}
By Theorem \ref{thm:successor-structure},
\[
\mathrm{Val}(c\oplus_S d)=\mathrm{Val}(c)+\mathrm{Val}(d).
\]
Since \(\mathrm{Val}\) is bijective, \(c\oplus_S d=c\oplus d\). The folding formula is exactly Corollary \ref{cor:add-as-fold}.
\end{proof}

\begin{theorem}[Multiplication generated by iterated successor]\label{thm:mul-from-successor}
Define
\[
c\otimes_S d:=\underbrace{c\oplus_S c\oplus_S\cdots\oplus_S c}_{\mathrm{Val}(d)\text{ times}},
\qquad
\text{and take }0_X\text{ if }\mathrm{Val}(d)=0.
\]
Then \(c\otimes_S d=c\otimes d\).
\leanverified{stableMul\_from\_successor}
\end{theorem}

\begin{proof}
By Corollary \ref{cor:add-from-successor}, we have \(\oplus_S=\oplus\). Hence
\[
\mathrm{Val}(c\otimes_S d)=\mathrm{Val}(c)\mathrm{Val}(d)=\mathrm{Val}(c\otimes d),
\]
and bijectivity of \(\mathrm{Val}\) gives the claim.
\end{proof}

\begin{proposition}[Minimality of successor among unary generators]\label{prop:successor-unary-minimal}
Let \(T:X_\infty\to X_\infty\) be a unary map with \(T(0_X)\ne 0_X\), and suppose that there is an integer \(r\ge 1\) such that
\[
\mathrm{Val}(T(c))=\mathrm{Val}(c)+r
\qquad(c\in X_\infty).
\]
If addition is to be recovered by the first-level iteration rule
\[
c\oplus d=T^{\circ\mathrm{Val}(d)}(c)
\qquad(c,d\in X_\infty),
\]
then \(r=1\) and \(T=S\). Thus the successor map is the unique nontrivial constant-step unary generator compatible with the iterated-readout formula for all stable addresses.
\end{proposition}

\begin{proof}
Taking values in the displayed iteration rule and using the assumed constant-step property gives
\[
\mathrm{Val}(c)+\mathrm{Val}(d)
=\mathrm{Val}(c\oplus d)
=\mathrm{Val}\bigl(T^{\circ\mathrm{Val}(d)}(c)\bigr)
=\mathrm{Val}(c)+r\,\mathrm{Val}(d)
\]
for all \(c,d\in X_\infty\). Choosing \(d=Z(1)\) gives \(r=1\). Hence \(\mathrm{Val}(T(c))=\mathrm{Val}(c)+1=\mathrm{Val}(S(c))\) for every \(c\), and bijectivity of \(\mathrm{Val}\) implies \(T(c)=S(c)\) for every \(c\). The assumption \(T(0_X)\ne 0_X\) only excludes the zero-step identity map, which cannot recover addition except when \(d=0_X\).
\end{proof}

\begin{corollary}[Sufficient primitive set and unary minimality]\label{cor:primitive-minimal-successor}
Arithmetic on stable addresses is generated by the successor map \(S\) together with folding \(\Fold_\infty\): addition is the first-level iterated readout of \(S\), while multiplication is the second-level iterated accumulation of that same readout. Within the class of nontrivial constant-step unary generators satisfying a first-level iterated-readout formula for all inputs, Proposition \ref{prop:successor-unary-minimal} shows that \(S\) is forced.
\end{corollary}

\begin{proof}
Generation is Corollary \ref{cor:add-from-successor} together with Theorem \ref{thm:mul-from-successor}. The asserted lower bound in the stated unary-generator interface is exactly Proposition \ref{prop:successor-unary-minimal}.
\end{proof}

\subsubsection{Two-layer readout from the single primitive \texorpdfstring{$\circ$}{o}}\label{subsec:arith-composition}
We now give a theorem-level formulation aligned with the chain
\[
\text{scan/iterate}\Longrightarrow \text{stable encoding}\Longrightarrow \text{normalized readout}.
\]
Instead of treating addition and multiplication as primitive, we retain only composition \(\circ\) as the most basic operation and return to canonical stable addresses by folding.

Let \(A\) be a nonempty set and let \(g:A\to A\) be an endomorphism whose iterates \(g^{\circ n}\) are pairwise distinct. Write
\[
S_0:=\langle g\rangle=\{g^{\circ n}:n\in\NN\}\subseteq \mathrm{End}(A),
\qquad
S_1:=\{I_n:n\in\NN\}\subseteq \mathrm{End}(\mathrm{End}(A)),
\]
where \(I_n(f):=f^{\circ n}\). The only operation on either level is composition \(\circ\).

\begin{theorem}[A single primitive \texorpdfstring{$\circ$}{o} generates both \texorpdfstring{$+$}{+} and \texorpdfstring{$\times$}{x} through two-layer readout]\label{thm:composition-two-layer}
The maps
\[
\iota_0:(\NN,+)\longrightarrow (S_0,\circ),
\iota_0(n):=g^{\circ n},
\qquad
g^{\circ 0}:=\mathrm{id}_A,
\]
and
\[
\iota_1:(\NN,\times)\longrightarrow (S_1,\circ),
\qquad
\iota_1(n):=I_n,
\]
are injective homomorphisms. In particular, for all \(m,n\in\NN\),
\[
\iota_0(m+n)=\iota_0(m)\circ \iota_0(n),
\qquad
\iota_1(mn)=\iota_1(m)\circ \iota_1(n).
\]
\end{theorem}

\begin{proof}
By the definition of iteration and associativity of composition,
\[
g^{\circ(m+n)}=g^{\circ m}\circ g^{\circ n},
\]
so \(\iota_0(m+n)=\iota_0(m)\circ \iota_0(n)\). On the other hand, for any \(f\in\mathrm{End}(A)\),
\[
(\iota_1(m)\circ \iota_1(n))(f)
=I_m(I_n(f))
=I_m(f^{\circ n})
=(f^{\circ n})^{\circ m}
=f^{\circ(nm)}
=I_{mn}(f).
\]
Since the two endomorphisms agree on every \(f\), we conclude that \(\iota_1(m)\circ \iota_1(n)=\iota_1(mn)\). Injectivity of \(\iota_0\) is the standing hypothesis. If \(\iota_1(m)=\iota_1(n)\), then applying both sides to \(g\) gives \(g^{\circ m}=g^{\circ n}\), hence \(m=n\). Thus \(\iota_1\) is injective as well.
\end{proof}

\begin{theorem}[One-layer obstruction for multiplicative readout]\label{thm:one-layer-composition-obstruction}
Let \(A\) be a nonempty set and let \(h:A\to A\) have pairwise distinct iterates. Suppose a readout of stable addresses into the one-generated composition semigroup \(\langle h\rangle\) has the form
\[
\Psi(c):=h^{\circ\mathrm{Val}(c)}.
\]
Then composition in \(\langle h\rangle\) realizes stable addition:
\[
\Psi(c\oplus d)=\Psi(c)\circ\Psi(d)
\qquad(c,d\in X_\infty),
\]
but it cannot realize stable multiplication by the same one-layer rule for all inputs. More precisely, there do not exist pairwise distinct iterates of \(h\) for which
\[
\Psi(c\otimes d)=\Psi(c)\circ\Psi(d)
\qquad(c,d\in X_\infty)
\]
holds identically. Thus the second layer in Theorem \ref{thm:composition-two-layer} is forced in the cyclic-iterate readout interface.
\end{theorem}

\begin{proof}
The additive identity follows from Theorem \ref{thm:add-definitional} and the law for iterates:
\[
\Psi(c\oplus d)
=h^{\circ\mathrm{Val}(c\oplus d)}
=h^{\circ(\mathrm{Val}(c)+\mathrm{Val}(d))}
=h^{\circ\mathrm{Val}(c)}\circ h^{\circ\mathrm{Val}(d)}.
\]
Assume, for contradiction, that the same one-layer composition rule realizes multiplication. Since the iterates of \(h\) are pairwise distinct, equality
\[
h^{\circ\mathrm{Val}(c\otimes d)}=h^{\circ(\mathrm{Val}(c)+\mathrm{Val}(d))}
\]
forces
\[
\mathrm{Val}(c)\mathrm{Val}(d)=\mathrm{Val}(c)+\mathrm{Val}(d)
\]
for all \(c,d\in X_\infty\). Taking \(c=Z(2)\) and \(d=Z(3)\) gives \(6=5\), a contradiction. Therefore multiplicative readout cannot be realized on the same cyclic-iterate layer; the endomorphism-of-endomorphisms layer \(I_n(f)=f^{\circ n}\) is precisely what converts composition into multiplication of exponents.
\end{proof}

\begin{theorem}[Two-layer composition realization of the stable semiring]\label{thm:arith-composition}
Retain Definitions \ref{def:stable-add} and \ref{def:stable-mul}, and let \(\mathrm{Val}:X_\infty\to\NN\) be the value map. For any pair \((A,g)\) as in Theorem \ref{thm:composition-two-layer}, define the process-semantics maps
\[
\Phi_0:X_\infty\to \mathrm{End}(A),
\qquad
\Phi_0(c):=g^{\circ\mathrm{Val}(c)},
\]
\[
\Phi_1:X_\infty\to \mathrm{End}(\mathrm{End}(A)),
\qquad
\Phi_1(c):=I_{\mathrm{Val}(c)}.
\]
Then for all \(c,d\in X_\infty\),
\[
\Phi_0(c\oplus d)=\Phi_0(c)\circ\Phi_0(d),
\qquad
\Phi_1(c\otimes d)=\Phi_1(c)\circ\Phi_1(d).
\]
That is, both \(\oplus\) and \(\otimes\) on stable addresses are realized by composition, but on two different layers of objects; Theorem \ref{thm:one-layer-composition-obstruction} shows that this layer separation is necessary for cyclic-iterate readouts.
\end{theorem}

\begin{proof}
By the definition of \(\oplus\) and the first identity in Theorem \ref{thm:composition-two-layer},
\[
\Phi_0(c\oplus d)
=g^{\circ\mathrm{Val}(c\oplus d)}
=g^{\circ(\mathrm{Val}(c)+\mathrm{Val}(d))}
=g^{\circ\mathrm{Val}(c)}\circ g^{\circ\mathrm{Val}(d)}
=\Phi_0(c)\circ\Phi_0(d).
\]
Similarly, by the definition of \(\otimes\) and the second identity in Theorem \ref{thm:composition-two-layer},
\[
\Phi_1(c\otimes d)
=I_{\mathrm{Val}(c\otimes d)}
=I_{\mathrm{Val}(c)\mathrm{Val}(d)}
=I_{\mathrm{Val}(c)}\circ I_{\mathrm{Val}(d)}
=\Phi_1(c)\circ\Phi_1(d).
\]
\end{proof}

\begin{corollary}[Readout by composition lands back in stable address space]\label{cor:composition-pullback}
Let \(Z:=\mathrm{Val}^{-1}\). Since \(\iota_0\) and \(\iota_1\) are injective, the assignments
\[
\rho_0:S_0\to X_\infty,
\qquad
\rho_0(g^{\circ n}):=Z(n),
\]
\[
\rho_1:S_1\to X_\infty,
\qquad
\rho_1(I_n):=Z(n).
\]
Then for any \(c,d\in X_\infty\),
\[
c\oplus d=\rho_0\bigl(\Phi_0(c)\circ \Phi_0(d)\bigr),
\qquad
c\otimes d=\rho_1\bigl(\Phi_1(c)\circ \Phi_1(d)\bigr).
\]
\end{corollary}

\begin{proof}
Write \(c=Z(m)\) and \(d=Z(n)\). Then \(\Phi_0(c)=g^{\circ m}\), \(\Phi_0(d)=g^{\circ n}\), \(\Phi_1(c)=I_m\), and \(\Phi_1(d)=I_n\). Theorem \ref{thm:composition-two-layer} therefore gives
\[
\rho_0\bigl(\Phi_0(c)\circ \Phi_0(d)\bigr)=\rho_0(g^{\circ m}\circ g^{\circ n})=Z(m+n),
\qquad
\rho_1\bigl(\Phi_1(c)\circ \Phi_1(d)\bigr)=\rho_1(I_m\circ I_n)=Z(mn),
\]
and these identities are exactly the definitions of \(\oplus\) and \(\otimes\) transported back to \(X_\infty\).
\end{proof}

\begin{corollary}[Primitive compression with a sharp layer obstruction]\label{cor:primitive-free-arith}
At the level of configurations, arithmetic may be compressed to the following minimal generating data:
\begin{itemize}
  \item pointwise superposition of counting configurations, i.e. the free commutative monoid structure;
  \item the local Fibonacci rewriting relations, viewed up to congruence;
  \item the canonical-section map \(\Fold_\infty\), or equivalently the pair \((\mathrm{Val},Z)\).
\end{itemize}
At the process level, the same content may be expressed using composition together with \(\Fold_\infty\): addition and multiplication are the readouts of composition on the first and second layers respectively, followed by normalization back to \(X_\infty\). The layer distinction is not cosmetic: within the cyclic-iterate readout interface, Theorem \ref{thm:one-layer-composition-obstruction} rules out a single-layer realization of multiplication.
\end{corollary}

\begin{proof}
At the configuration level, the three ingredients are exactly those isolated by Theorem \ref{thm:monoid-quotient-is-N} and Corollary \ref{cor:fold-as-section}. At the process level, the two-layer readout statements are Theorems \ref{thm:composition-two-layer} and \ref{thm:arith-composition}, while the interface with addition is supplied by Corollary \ref{cor:add-as-fold}. The asserted obstruction is Theorem \ref{thm:one-layer-composition-obstruction}.
\end{proof}

\begin{remark}[Euclidean division as a choice of normal form]\label{rem:division-as-selection}
Given \(a,b\in X_\infty\) with \(\mathrm{Val}(b)\ge 1\), the integer-level quotient-and-remainder relation
\[
\mathrm{Val}(a)=\mathrm{Val}(q)\,\mathrm{Val}(b)+\mathrm{Val}(r),
\qquad
0\le \mathrm{Val}(r)<\mathrm{Val}(b),
\]
may be rewritten equivalently as the stable-address equation
\[
a=q\otimes b\oplus r,
\]
supplemented by the inequality constraint selecting the unique normal representative \((q,r)\). Thus division with remainder is not a new primitive operation; it is a problem of solving an equation inside the same rewriting groupoid and then choosing the canonical normal form. Its computability is inherited from the computability of \(\Fold_\infty\).
\end{remark}

\begin{proposition}[Computable realization]\label{prop:zeckendorf-arith-computable}
For any \(c,d\in X_\infty\), both \(c\oplus d\) and \(c\otimes d\) are computable in finitely many steps: one first performs ordinary integer addition or multiplication at the value level to obtain \(\mathrm{Val}(c)+\mathrm{Val}(d)\) or \(\mathrm{Val}(c)\mathrm{Val}(d)\), and then applies the greedy Zeckendorf algorithm to recover the unique normal form.
\end{proposition}

\begin{proof}
The values \(\mathrm{Val}(c)\) and \(\mathrm{Val}(d)\) are finite integers because \(c\) and \(d\) have finite support. Integer addition and multiplication are algorithmic, and the greedy Zeckendorf procedure returns the unique admissible representative of any nonnegative integer \cite{Zeckendorf1972,Koshy2018}. Hence both operations are computable.
\end{proof}
